% !TeX program = xelatex
\documentclass{cumcmthesis}
\usepackage{pdflscape}
\usepackage{makecell}
\usepackage{diagbox}
\usepackage{tikz}
\usepackage{fvextra}
\usepackage{multicol}
\usetikzlibrary{arrows.meta,positioning,shapes.geometric}
\graphicspath{{../figures/PDF/}}

\newcommand{\vct}[1]{\vec{#1}}
\newcommand{\dd}{\,\mathrm d}
\newcommand{\Tset}{\mathcal T}
\newcommand{\Qset}{\mathcal Q}
\newcommand{\Sset}{\mathcal S}
\newcommand{\clip}{\operatorname{clip}}
\newcommand{\meas}{\mu}
\newcommand{\Hs}{H_{\Sigma}}
\newcommand{\Hu}{H_{\cup}}
\emergencystretch=2em
\raggedbottom
\setlength{\textfloatsep}{8pt plus 2pt minus 2pt}
\setlength{\floatsep}{7pt plus 2pt minus 2pt}
\setlength{\intextsep}{7pt plus 2pt minus 2pt}
\newcommand{\steptext}[3]{\par\noindent\textbf{STEP#1：#2。}#3\par\vspace{0.10em}}
% 重点文字统一使用黑体/无衬线粗体，与正文宋体形成清晰视觉层级。
% 数学重点同时启用 boldmath，避免中文已加粗而数字/符号仍保持常规字重。
\renewcommand{\keyresult}[1]{{\sffamily\bfseries\boldmath #1}}
\newcommand{\inlinebox}[1]{\(\boxed{#1}\)}
\newcommand{\abslead}[1]{{\sffamily\bfseries #1}}
\newcommand{\absstrong}[1]{{\sffamily\bfseries\boldmath #1}}

\title{基于有限视线段完整遮蔽判据的烟幕干扰弹协同投放优化}

\begin{document}

\maketitle
\begin{abstract}
\setlength{\parindent}{2em}
\setlength{\parskip}{0.10em}
针对多架无人机投放烟幕干扰弹、持续遮蔽来袭导弹视线的任务规划问题，本文以真实目标的有限几何尺度为出发点，建立“\absstrong{有限视线段完整遮蔽判定—多烟幕联合覆盖—分层连续优化与严格复算}”的统一模型链。真实目标始终保持为半径 $7\,\mathrm m$、高 $10\,\mathrm m$ 的有限闭圆柱；对每一时刻，从导弹到目标可见点构造有限视线段集合，并以最坏视线到有效烟幕域的距离作为完整遮蔽判据，从而避免点目标和无限直线近似造成的误判。

\abslead{针对问题一，}在题给固定策略下建立有限圆柱首次可见点与点到有限线段距离模型，得到唯一有效遮蔽区间 $[8.056430,\,9.448087]\,\mathrm s$，\absstrong{有效遮蔽时长为 $1.391656494\,\mathrm s$}。连续判据与有限射线数值评价器分开表述；中、高两级离散结果仅差 $6.10\times10^{-6}\,\mathrm s$，独立有限圆柱几何核的最大距离差为 $8.66\times10^{-5}\,\mathrm m$。

\abslead{针对问题二至问题四，}依次释放单弹参数、同机三弹时序和三机独立航迹，将决策维数由4维扩展至8维和12维，但完整遮蔽几何核保持不变。问题二得到航向 $5.134472^\circ$、速度 $139.988137\,\mathrm{m/s}$ 的单弹方案，遮蔽时长为 \absstrong{$4.588059082\,\mathrm s$}，较固定策略增加 $3.196403\,\mathrm s$；问题三的三弹联合区间为 $[0.999999,\,7.403160]\,\mathrm s$，时长为 \absstrong{$6.403160400\,\mathrm s$}，其中第三烟幕槽在最终方案中的逐弹边际贡献为0；问题四由 FY1--FY3 形成三个互不重叠的有效区间，累计遮蔽 \absstrong{$11.669261475\,\mathrm s$}，说明主要协同机制是前、中、后三段的时域接力。

\abslead{针对问题五，}建立五机、每机三个连续烟幕槽的40维联合模型，并证明在“烟幕之间无负向作用”及本题任务时域下，该固定槽位表示与“每机至多三枚”的最优目标值等价。模型不设置 UAV--导弹排他分配变量，所有有效烟幕均可同时参与三枚导弹的几何评价；15个单机--单导弹局部子问题仅用于构造结构化初值。全局搜索记录中，结构化/混合初始化支路的严格候选最高达到 $18.93\sim19.18\,\text{导弹}\cdot\mathrm s$，纯随机冷启动支路最高仅为 $1.44\,\text{导弹}\cdot\mathrm s$，随后经五机变量块精修与时序抛光得到最终方案。三枚导弹的有效遮蔽时长分别为 $10.083504639$、$6.718437500$ 和 $7.362644043\,\mathrm s$，主效能为 \absstrong{$H_\Sigma=24.164586182\,\text{导弹}\cdot\mathrm s$}，墙钟并集为 $17.418010254\,\mathrm s$。

全题最终结果均经过分辨率加密与独立有限圆柱几何核复核；问题五 L2--L3 的最大分导弹时长差为 $1.094\times10^{-3}\,\mathrm s$，全题独立几何核最大差不超过 $3.67\times10^{-4}\,\mathrm m$。因此本文将高维启发式结果表述为经多层证据支持的\absstrong{稳定高质量方案}，不作理论全局最优声明。
\end{abstract}

\keywords{烟幕干扰；有限视线段；完整遮蔽；差分进化；协同优化}
\clearpage

\section{问题重述}

\subsection{任务背景与统一约束}
题目给出三枚来袭导弹、五架无人机及一个圆柱形真实目标。导弹均以 $300\,\mathrm{m/s}$ 的恒定速率飞向位于坐标原点的假目标；无人机受领任务后可瞬时确定水平航向和飞行速度，此后保持高度并作匀速直线运动，速度范围为 $70\sim140\,\mathrm{m/s}$。无人机投放的烟幕弹在重力作用下运动，起爆后形成有效半径 $10\,\mathrm m$ 的球状烟幕云团，云团在 $20\,\mathrm s$ 有效期内以 $3\,\mathrm{m/s}$ 匀速下沉。同一无人机相邻两次投弹至少间隔 $1\,\mathrm s$。上述条件均直接来自题面\cite{ref1}。

真实目标不是空间点，而是半径 $7\,\mathrm m$、高度 $10\,\mathrm m$ 的有限圆柱。因此，各问的共同前提是：烟幕是否真正阻断导弹到\emph{整个有限目标可见部分}的视线，而不是仅遮住目标中心或与某条无限直线相交。

\subsection{五个问题的任务链}
题目五问沿“固定策略评价--单弹优化--同机多弹--多机协同--多导弹联合”的方向逐级扩展：
\begin{enumerate}[label=\arabic*.]
  \item 在 FY1 的飞行方向、速度、投弹时刻和起爆延迟均给定时，计算单枚烟幕弹对 M1 的有效遮蔽时长；
  \item 释放 FY1 的航向、速度、投弹时刻和起爆延迟，求单机单弹的最优投放策略；
  \item FY1 连续投放三枚烟幕弹，在共享航向、共享速度及相邻投弹间隔约束下设计三弹时序；
  \item FY1--FY3 各投放一枚烟幕弹，联合确定三架无人机的航迹与起爆参数，使对 M1 的有效遮蔽尽可能延长；
  \item 五架无人机各至多投放三枚烟幕弹，对 M1--M3 实施联合干扰，并给出完整可执行的投放方案及三枚导弹的有效遮蔽结果。
\end{enumerate}

\section{问题分析}

五个问题共用同一组运动规律和目标几何，但数学难点逐问增加。全题首先需要建立一个能在有限圆柱目标上统一使用的\textbf{完整遮蔽评价器}；随后只改变决策变量的组织方式和协同层级。由此可将全题组织为“公共运动关系--有限目标几何判定--单烟幕参数优化--多烟幕/多平台协同--多导弹联合评价”的递进链。

\subsection{问题一：从有限目标构造完整遮蔽判据}
本问没有优化变量，关键不是运动轨迹本身，而是“有效遮蔽”的几何定义。对任一时刻，导弹到目标可见点形成一束\emph{有限视线段}；单枚烟幕只有同时截获这束视线中的全部成员，才能形成完整遮蔽。因此需要先定义点到有限线段的距离，再对有限圆柱的可见视线集合取最不利值。该评价器一旦建立，可直接作为后续四问的统一几何核。

\subsection{问题二：四维参数化下的单机单弹优化}
FY1 一旦确定航向和速度就不再调整，投放点与起爆点也分别由投弹时刻和起爆延迟唯一决定，因此无需把空间坐标重复作为独立决策变量。问题二可压缩为“航向--速度--投弹时刻--起爆延迟”四个连续变量。其目标值由完整遮蔽时间集合的长度决定，随参数变化会出现区间生成、消失和边界切换，因而属于非凸、非光滑的连续优化问题。

\subsection{问题三：共享平台约束下的三弹联合遮蔽}
三枚烟幕共享同一架无人机的航向和速度，同时受投弹间隔约束。多烟幕协同不能简单把三枚烟幕的独立遮蔽时长相加：在同一时刻，不同烟幕可以分别遮挡不同的目标视线，因此应对每条视线先选择最有利的有效烟幕，再对全部可见视线取最坏情形。该结构把“同机时序约束”和“空间互补遮蔽”耦合为一个八维连续问题。

\subsection{问题四：不同初始位置下的三机时空协同}
三架无人机的航向和速度互不共享，但初始位置差异显著。与问题三相比，本问解除同机三弹的共享平台约束，却增加三组独立航迹变量。需要判断三架无人机的有效区间是依靠同时空间拼接扩展，还是在不同时间段形成接力。因而除了求最大遮蔽时长，还应分析各机独立区间与联合区间之间的集合关系。

\subsection{问题五：五机多弹对三导弹的联合决策}
问题五同时包含五组固定航迹、最多十五个投弹事件和三枚来袭导弹。由于任一有效烟幕都可能依据实时几何关系帮助任一导弹，若预先设置排他的“无人机--导弹”分配变量，反而可能截断真实协同。因此正式模型应保留全部烟幕对三枚导弹的自由几何贡献，并以三枚导弹各自有效时间的总量作为主效能指标；同时另行统计墙钟并集，以解释双重、三重同时干扰产生的协同增益。

综上，全题技术路线如图\ref{fig:route}所示。图中公共几何判据从问题一向后继问题持续传递，而问题三、问题四分别提供“同机多弹”和“多机协同”的结构基础，最终汇入问题五。

\begin{figure}[tbp]
  \centering
  \includegraphics[width=0.95\textwidth]{F00_全题技术路线图.pdf}
  \caption{全题问题链与技术路线}
  \label{fig:route}
\end{figure}

\section{模型假设}

题面已明确的导弹速度、无人机匀速等高直线飞行、烟幕有效半径与下沉速度属于已知事实，不再重复列为模型假设。为补足题面未规定但计算所必需的物理细节，作如下简化：
\begin{enumerate}[label=\arabic*.]
  \item \textbf{烟幕弹自由飞行阶段忽略空气阻力与风场。} 烟幕弹脱离无人机瞬间继承载机的水平速度，竖直初速度取零；该假设使投放到起爆阶段可用经典抛体运动描述。若存在强风或显著气动阻力，起爆点会产生额外水平/竖直偏移。
  \item \textbf{起爆后的有效云团采用硬边界球模型。} 在题给 $20\,\mathrm s$ 有效期内，以云团中心 $10\,\mathrm m$ 内为有效遮蔽区域；除题给 $3\,\mathrm{m/s}$ 竖直下沉外，不再加入水平漂移、半径扩张和浓度梯度。该简化与题面给出的有效范围直接对应，但不能描述真实烟幕的非均匀扩散。
  \item \textbf{多枚烟幕之间不存在负向相互作用。} 不考虑云团碰撞、稀释或遮蔽能力相互削弱；多烟幕仅通过几何覆盖形成联合效应。因而增加一枚可行烟幕不会使已有视线的几何遮蔽能力变差。
  \item \textbf{飞行与起爆过程按确定性计划执行。} 不考虑导航、投放与引信控制误差，题面坐标和时刻视为确定量。该假设用于求题设条件下的基准最优策略，实际部署时可在本文模型外进一步加入执行误差与鲁棒裕度。
\end{enumerate}

\section{主要符号说明}
\begingroup
\normalsize
\renewcommand{\arraystretch}{1.22}
\setlength{\tabcolsep}{6pt}
\begin{center}
\begin{tabularx}{\textwidth}{>{\centering\arraybackslash}p{0.17\textwidth} X >{\centering\arraybackslash}p{0.18\textwidth}}
\toprule
\textbf{符号} & \textbf{含义} & \textbf{单位} \\
\midrule
$t$ & 连续时间 & s \\
$\vec m_j(t)$ & 导弹 $M_j$ 在时刻 $t$ 的位置向量 & m \\
$\vec f_i(t)$ & 无人机 FY$i$ 在时刻 $t$ 的位置向量 & m \\
$\theta_i$ & FY$i$ 的水平航向角 & rad 或 $^\circ$ \\
$v_i$ & FY$i$ 的飞行速度 & m/s \\
$t_{ik}$ & FY$i$ 第 $k$ 枚烟幕弹的投放时刻 & s \\
$\delta_{ik}$ & FY$i$ 第 $k$ 枚烟幕弹的起爆延迟 & s \\
$\tau_{ik}$ & FY$i$ 第 $k$ 枚烟幕弹的起爆时刻 & s \\
$\vec c_{ik}(t)$ & 第 $(i,k)$ 枚有效烟幕云团中心位置 & m \\
$\mathcal T$ & 真实目标对应的有限闭圆柱集合 & --- \\
$\mathcal V_j(t)$ & 从 $M_j$ 观察到的目标可见点集合 & --- \\
$d_{\rm seg}$ & 点到有限视线段的最短欧氏距离 & m \\
$D_j(t)$ & 对 $M_j$ 的完整遮蔽最坏距离 & m \\
$E_j$ & $M_j$ 的有效遮蔽时间集合 & --- \\
$H_j$ & $M_j$ 的有效遮蔽时长 $\mu(E_j)$ & s \\
$H_\Sigma$ & 三枚导弹有效遮蔽时长之和 & 导弹$\cdot$s \\
$H_\cup$ & 三枚导弹有效时间集合的墙钟并集长度 & s \\
\bottomrule
\end{tabularx}
\vspace{0.12em}
\begin{minipage}{\textwidth}
{\small\color{cumcmgray}\raggedright 注：仅列出跨问反复出现且评委需要回查的主要符号。局部投影参数、算法迭代索引和临时变量均在首次出现处定义。$H_j$ 始终按导弹编号索引，不再用 $H_1,H_2,H_3,H_4$ 表示问题编号结果。\par}
\end{minipage}
\end{center}
\endgroup

\section{模型准备}

本章只建立五问共同使用的坐标、运动关系与真实目标几何。完整遮蔽判据在问题一中建立，随后通过统一评价函数自然推广到多烟幕和多导弹情形。

\subsection{导弹与无人机运动}
以假目标为坐标原点 $O=(0,0,0)$。第 $j$ 枚导弹初始位置为 $\vec m_j(0)$，其单位飞行方向为
\begin{equation}
\vec e_j=-\frac{\vec m_j(0)}{\lVert\vec m_j(0)\rVert},
\qquad
\vec m_j(t)=\vec m_j(0)+300t\vec e_j.
\label{eq:missile}
\end{equation}
对应到达假目标的时刻为 $T_j=\lVert\vec m_j(0)\rVert/300$。对多导弹联合问题，统一记任务时间上界为 $T_{\max}=\max_j T_j$；问题一至问题四只需使用 $T_1$。

设 FY$i$ 的初始位置为 $\vec f_i(0)=(x_i,y_i,h_i)$，水平航向角为 $\theta_i$，速度为 $v_i$。定义水平单位方向
\begin{equation}
\vec u_i=(\cos\theta_i,\sin\theta_i,0),
\qquad
\vec f_i(t)=\vec f_i(0)+v_it\vec u_i,
\quad 70\le v_i\le140.
\label{eq:uav}
\end{equation}
初始坐标均采用题面数据，后续各问只引用对象编号，不重复列式。

\subsection{烟幕弹与云团运动}
FY$i$ 的第 $k$ 枚烟幕弹在 $t_{ik}$ 投放，延迟 $\delta_{ik}$ 起爆，故起爆时刻为
\begin{equation}
\tau_{ik}=t_{ik}+\delta_{ik}.
\label{eq:burst-time}
\end{equation}
在假设1下，烟幕弹继承无人机水平速度并作抛体运动。记 $\vec e_z=(0,0,1)$，则起爆点为
\begin{equation}
\vec b_{ik}=\vec f_i(0)+v_i\tau_{ik}\vec u_i
-\frac12g\delta_{ik}^{2}\vec e_z,
\qquad g=9.8\,\mathrm{m/s^2}.
\label{eq:burst-point}
\end{equation}
起爆前不触地给出延迟上界
\begin{equation}
0\le\delta_{ik}\le\sqrt{\frac{2h_i}{g}}.
\label{eq:delay-cap}
\end{equation}
起爆后，云团中心在有效期内满足
\begin{equation}
\vec c_{ik}(t)=\vec b_{ik}-3(t-\tau_{ik})\vec e_z,
\qquad t\in[\tau_{ik},\tau_{ik}+20].
\label{eq:cloud}
\end{equation}
记时刻 $t$ 的有效烟幕中心集合为 $\mathcal C(t)$。该集合只包含已经起爆且仍处于 $20\,\mathrm s$ 有效期内的烟幕。

\subsection{真实目标几何}
真实目标严格保持为题目给出的有限闭圆柱
\begin{equation}
\mathcal T=\left\{(x,y,z):x^2+(y-200)^2\le49,\ 0\le z\le10\right\}.
\label{eq:target}
\end{equation}
其半径为 $7\,\mathrm m$、高度为 $10\,\mathrm m$。后文所有视线、距离和遮蔽条件均以该有限实体为终点，不将其退化为中心点。

\section{问题一：固定策略的完整遮蔽评价}

\subsection{固定策略与有限视线段距离}
问题一中 FY1 的控制参数均由题面给定。其水平航向指向假目标，故 $\theta=\pi$，速度为 $120\,\mathrm{m/s}$；投放时刻取 $t_0=1.5\,\mathrm s$，起爆延迟为 $\delta=3.6\,\mathrm s$，从而起爆时刻
\[
\tau=t_0+\delta=5.1\,\mathrm s.
\]
由式\eqref{eq:uav}--\eqref{eq:burst-point}得到投放点和起爆点
\[
\vec p=(17620,0,1800)\ \mathrm m,
\qquad
\vec b=(17188,0,1736.496)\ \mathrm m.
\]

任取时刻 $t$，导弹位置为 $\vec a=\vec m_1(t)$，目标上一可见点为 $\vec q$，烟幕中心为 $\vec c(t)$。导弹到该目标点的真实视线是有限线段
\[
[\vec a,\vec q]=\{\vec a+\lambda(\vec q-\vec a):0\le\lambda\le1\}.
\]
令
\begin{equation}
\lambda_0=
\frac{(\vec c-\vec a)\cdot(\vec q-\vec a)}{\lVert\vec q-\vec a\rVert^2},
\qquad
\lambda^*=\clip(\lambda_0,0,1),
\label{eq:lambda}
\end{equation}
则点到有限视线段的最短距离为
\begin{equation}
 d_{\rm seg}\!\left(\vec c,[\vec a,\vec q]\right)
 =\left\lVert\vec c-
 \bigl[\vec a+\lambda^*(\vec q-\vec a)\bigr]\right\rVert .
\label{eq:dseg}
\end{equation}
式\eqref{eq:lambda}中的截断是本题几何判定的重要环节：当正交投影落在线段外时，最近点自动退化为导弹端点或目标端点，因此不会把“烟幕位于导弹后方或目标后方但与无限直线很近”的情况误判为有效遮蔽。几何构造见图\ref{fig:q1-geometry}。

\begin{figure}[tbp]
  \centering
  \includegraphics[width=0.94\textwidth]{F01_Q1有限视线段几何机理.pdf}
  \caption{有限视线段完整遮蔽判据的几何构造}
  \label{fig:q1-geometry}
\end{figure}

\subsection{有限圆柱可见集与完整遮蔽判据}
从导弹位置 $\vec m_1(t)$ 向目标 $\mathcal T$ 发出所有能够击中圆柱的射线，并在每条射线上只保留与 $\mathcal T$ 的\textbf{第一个}交点。由这些首次交点构成目标在时刻 $t$ 的可见点集合 $\mathcal V_1(t)$。于是，对整个有限目标的最坏视线距离定义为
\begin{equation}
D_1(t)=
\max_{\vec q\in\mathcal V_1(t)}
 d_{\rm seg}\!\left(
 \vec c(t),[\vec m_1(t),\vec q]
 \right).
\label{eq:q1-D}
\end{equation}
烟幕有效半径为 $R_s=10\,\mathrm m$，故完整遮蔽条件为
\begin{equation}
D_1(t)\le R_s.
\label{eq:q1-block}
\end{equation}

\begin{propositionbox}
\textbf{命题1（完整遮蔽等价判据）}\quad
在球形有效烟幕域 $B(\vec c(t),R_s)$ 下，式\eqref{eq:q1-block}成立，当且仅当从 $M_1$ 指向真实目标可见部分的每一条有限视线段都与该烟幕域相交。

\textbf{证明：}
若 $D_1(t)\le R_s$，由最大值定义，对任意 $\vec q\in\mathcal V_1(t)$ 均有
$d_{\rm seg}(\vec c,[\vec m_1,\vec q])\le R_s$，故相应线段与烟幕球相交。反之，若所有可见线段均与烟幕球相交，则每条线段到球心的距离都不超过 $R_s$，取最大仍有 $D_1(t)\le R_s$。证毕。\proofend
\end{propositionbox}

因此有效时间集合与时长统一写为
\begin{equation}
E_1=\{t\in[\tau,\tau+20]\cap[0,T_1]:D_1(t)\le10\},
\qquad
H_1=\mu(E_1).
\label{eq:q1-set}
\end{equation}
这一记号在后续各问保持不变：问题一至问题四均针对 $M_1$，因此最终结果始终是不同策略下的 $H_1$，不再以 $H_2,H_3,H_4$ 表示问题编号。

\subsection{有限圆柱首次交点与数值评价器}
式\eqref{eq:q1-D}是连续可见集上的数学定义。为数值求值，需要构造有限圆柱的可见射线。记圆柱几何中心为 $\vec o=(0,200,5)$，外接球半径为
\[
R_B=\sqrt{7^2+5^2}=\sqrt{74}.
\]
在通过 $\vec o$ 且垂直于视轴 $\vec o-\vec m_1(t)$ 的法平面内，用外接球切锥确定需要扫描的二维圆盘。若 $d=\lVert\vec o-\vec m_1(t)\rVert$，相应圆盘半径为
\begin{equation}
R_{\Pi}(t)=\frac{dR_B}{\sqrt{d^2-R_B^2}}.
\label{eq:plane-radius}
\end{equation}

对圆盘中的任一方向 $\vec u$，射线写为
\[
\vec r(\rho)=\vec m_1(t)+\rho\vec u,
\qquad \rho\ge0.
\]
设 $\vec m_1(t)=(m_x,m_y,m_z)$、$\vec u=(u_x,u_y,u_z)$。侧面交点满足
\begin{equation}
A\rho^2+B\rho+C=0,
\label{eq:cyl-side}
\end{equation}
其中
\[
A=u_x^2+u_y^2,\quad
B=2[m_xu_x+(m_y-200)u_y],\quad
C=m_x^2+(m_y-200)^2-49.
\]
只保留 $\rho\ge0$ 且 $0\le m_z+\rho u_z\le10$ 的根。对上下底面 $z_b\in\{0,10\}$，若 $u_z\ne0$，候选根为
\begin{equation}
\rho_b=\frac{z_b-m_z}{u_z},
\label{eq:cyl-cap}
\end{equation}
并要求交点落在半径 $7\,\mathrm m$ 的圆盘内。在全部合法侧面根和底面根中取最小正根 $\rho_{\rm hit}$，即可得到该射线对应的首次可见点
\begin{equation}
\vec q=\vec m_1(t)+\rho_{\rm hit}\vec u.
\label{eq:first-hit}
\end{equation}

需要强调，\eqref{eq:q1-D}中的 $D_1(t)$ 是连续可见集上的精确定义，而程序只能在有限射线集合上计算其近似值。令 $\mathcal R(t;N)\subset\mathcal V_1(t)$ 表示分辨率参数为 $N$ 时得到的自适应首次交点集合，定义
\begin{equation}
\widehat D_1(t;N)=
\max_{\vec q\in\mathcal R(t;N)}
 d_{\rm seg}\!\left(\vec c(t),[\vec m_1(t),\vec q]\right).
\label{eq:q1-Dhat}
\end{equation}
本文不把有限采样的 $\widehat D_1$ 宣称为连续最大值的解析精确计算，而是通过\textbf{法平面逐级加密、最坏方向局部再加密以及独立有限圆柱几何核}检验其数值稳定性。这样将“数学判据的精确性”和“数值评价器的近似性”明确分开。

时间方向令 $g(t)=\widehat D_1(t;N)-10$。先粗扫描烟幕有效窗，定位 $g(t)$ 的符号变化区间，再对每个边界执行二分，并同步加密当前最坏视线邻域。求解过程概括为：
\steptext{1}{固定轨迹计算}{由题给参数得到导弹、无人机、烟幕弹及云团中心轨迹，并确定有效时间窗。}
\steptext{2}{有限目标可见集离散}{在视轴法平面生成射线，按式\eqref{eq:cyl-side}--\eqref{eq:cyl-cap}求有限圆柱首次交点。}
\steptext{3}{遮蔽边界初定位}{计算 $\widehat D_1(t;N)-10$，记录符号发生变化的时间小区间。}
\steptext{4}{时间--几何联合精化}{对候选边界二分时间，同时围绕当前最坏可见方向增加射线密度。}
\steptext{5}{最高精度复算}{合并有效子区间，使用最高分辨率和独立几何核重新计算边界与距离。}

\begin{center}
\begin{minipage}{0.94\textwidth}
\centering\textbf{算法1：固定策略的完整遮蔽边界计算}\par\vspace{0.18em}
\begin{tabularx}{\textwidth}{>{\raggedleft\arraybackslash}p{0.045\textwidth}X}
\toprule
\multicolumn{2}{l}{\textbf{输入：}固定飞行参数、投放/起爆时序，时间与法平面分辨率。}\\
\multicolumn{2}{l}{\textbf{输出：}有效集合 $E_1$、时长 $H_1$、数值收敛与独立几何复核误差。}\\
\midrule
1 & 计算 $\vec m_1(t)$、$\vec c(t)$ 以及烟幕有效时间窗；\\
2 & 对每个粗时间点生成视轴法平面射线，并求有限圆柱的首次交点集合 $\mathcal R(t;N)$；\\
3 & 计算 $\widehat D_1(t;N)$，记录与 $10\,\mathrm m$ 阈值发生状态切换的时间小区间；\\
4 & 对候选边界二分，并在当前最坏可见方向局部加密；\\
5 & 合并满足阈值的时间区间，以最高精度和独立几何核完成认证。\\
\bottomrule
\end{tabularx}
\end{minipage}
\end{center}

\subsection{结果与数值认证}
最高精度复算得到唯一有效遮蔽区间
\[
E_1=[8.056430,\ 9.448087]\ \mathrm s,
\]
从而
\begin{equation}
\boxed{H_1=1.391656494\ \mathrm s}.
\label{eq:q1-result}
\end{equation}
图\ref{fig:q1-result}把阈值边界、离散加密和独立几何核复核放在同一证据链中。进入与退出时刻均对应最坏视线距离达到 $10\,\mathrm m$；有效区间内部的最小距离约为 $4.735\,\mathrm m$，说明中段具有约 $5.265\,\mathrm m$ 的几何裕度。

\begin{figure}[tbp]
  \centering
  \includegraphics[width=0.93\textwidth]{F02_Q1阈值边界与数值认证.pdf}
  \caption{问题一完整遮蔽阈值、有效时间窗与数值认证}
  \label{fig:q1-result}
\end{figure}

三档离散设置得到的遮蔽时长分别约为 $1.391699\,\mathrm s$、$1.391650\,\mathrm s$ 和 $1.391656494\,\mathrm s$，中、高两级差仅 $6.10\times10^{-6}\,\mathrm s$；独立有限圆柱几何核在复核时刻上的最大绝对距离差约为 $8.66\times10^{-5}\,\mathrm m$。相较于 $10\,\mathrm m$ 的有效半径，这一几何差异处于 $10^{-5}$ 的相对量级。故后续各问均继续使用同一完整遮蔽评价器，并在最终结果处执行同等级别的严格复算。

\FloatBarrier

\section{问题二：单机单弹的四维连续优化}

本问完整继承问题一的导弹运动、烟幕运动、有限圆柱可见集和完整遮蔽评价器，仅把题给固定策略释放为四个连续决策变量。其“继承--调整--输出”关系见表\ref{tab:q2-inherit}。

\begin{table}[tbp]
\centering
\caption{问题二相对问题一的继承--调整--输出关系}
\label{tab:q2-inherit}
\small
\begin{tabularx}{\textwidth}{p{0.27\textwidth} X X}
\toprule
\textbf{继承} & \textbf{调整} & \textbf{输出} \\
\midrule
$\vec m_1(t)$、$\mathcal T$、$\mathcal V_1(t)$、$d_{\rm seg}$ 与完整遮蔽判据 & 将 FY1 航向、速度、投放时刻、起爆延迟由题给常量释放为连续变量，不把投放点和起爆点重复作为自由变量 & 最优四维策略、有效集合 $E_1$、遮蔽时长 $H_1$ 与搜索/几何稳定性证据 \\
\bottomrule
\end{tabularx}
\end{table}

\subsection{四维决策变量与完整模型}
问题二仍由 FY1 干扰 M1。由于无人机一旦确定航向和速度便保持匀速直线运动，投放点和起爆点均可由时序参数唯一确定，因此选取
\begin{equation}
\vec x=(\theta,v,t_0,\delta)\in\mathbb R^4
\label{eq:q2-x}
\end{equation}
作为最小充分决策向量，其中 $t_0$ 为投放时刻。给定 $\vec x$ 后，由第5章运动关系得到烟幕轨迹，并用问题一的连续完整遮蔽判据定义
\begin{equation}
E_1(\vec x)=\{t:D_1(t;\vec x)\le10\},
\qquad
H_1(\vec x)=\mu(E_1(\vec x)).
\label{eq:q2-obj}
\end{equation}
于是问题二的完整模型为
\begin{equation}
\left\{
\begin{aligned}
\max_{\theta,v,t_0,\delta}\quad & H_1(\vec x),\\
\text{s.t.}\quad
&0\le\theta<2\pi,\\
&70\le v\le140,\\
&0\le t_0\le T_1,\\
&0\le\delta\le\sqrt{2h_1/g},\\
&t_0+\delta\le T_1.
\end{aligned}
\right.
\label{eq:q2-model}
\end{equation}

四个变量存在明显耦合：$v$ 与 $\theta$ 决定水平可达方向，$t_0$ 决定释放前的飞行距离，而 $\delta$ 同时产生 $v\delta$ 的水平推进和 $g\delta^2/2$ 的高度损失。更重要的是，$H_1$ 由阈值事件的时间集合测度构成，参数的微小变化可能使有效区间产生、合并或消失，因此目标函数具有非光滑和多盆地特征。

\subsection{全局搜索、独立对照与盆地精修}
主搜索采用差分进化（Differential Evolution, DE）\cite{ref2,ref3,ref9}。设第 $g$ 代第 $i$ 个候选为 $\vec x_{i,g}$，当前最好个体为 $\vec x_{\rm best,g}$，主支路使用 best/1/bin 变异
\begin{equation}
\vec v_{i,g}=\vec x_{\rm best,g}
+F(\vec x_{r_1,g}-\vec x_{r_2,g}),
\label{eq:de-mutation}
\end{equation}
再通过二项交叉形成试验个体 $\vec u_{i,g}$。若试验个体满足全部物理约束且严格评价值不低于父代，则完成贪婪替换。由于 DE 属于随机启发式，本问并不依赖单次运行：3 次独立 DE 用于发现多个优质盆地，并以双重退火\cite{ref4}提供不同搜索机制的可行对照；所有候选最终都必须回到同一遮蔽评价器和同一精度下复算。

求解过程按以下五步执行：
\steptext{1}{全局探索}{在式\eqref{eq:q2-model}的四维可行域内执行3次独立 DE，同时生成双重退火对照候选。}
\steptext{2}{统一重排}{将各支路候选全部送入同一中高精度 $H_1(\vec x)$ 评价器，消除搜索阶段精度差异造成的排序偏差。}
\steptext{3}{多盆地保留}{综合目标值与参数空间距离保留4个互异优质盆地，避免局部精修只围绕一个随机初值展开。}
\steptext{4}{局部缩域}{在各盆地中逐轮缩小参数盒并继续无导数搜索，提升局部边界定位精度。}
\steptext{5}{严格认证}{对盆地最优候选重新定位遮蔽边界，并执行时间--视线分辨率加密与独立几何核复核。}

\begin{center}
\begin{minipage}{0.94\textwidth}
\centering\textbf{算法2：单机单弹的多盆地连续优化}\par\vspace{0.18em}
\begin{tabularx}{\textwidth}{>{\raggedleft\arraybackslash}p{0.045\textwidth}X}
\toprule
\multicolumn{2}{l}{\textbf{输入：}四维可行域 $\Omega$、完整遮蔽评价器 $H_1(\vec x)$、全局与局部搜索预算。}\\
\multicolumn{2}{l}{\textbf{输出：}FY1 航向、速度、投放/起爆时序与最高精度遮蔽时长。}\\
\midrule
1 & 独立初始化多组 DE 种群，并运行双重退火生成异质候选；\\
2 & 对每个试验个体执行可行域修复，调用统一几何评价器计算 $H_1$；\\
3 & 汇总前列候选并统一严格重排，按目标值和参数距离提取互异盆地；\\
4 & 对各盆地执行局部缩域精修；\\
5 & 对最优候选执行分层数值加密、边界重定位和独立有限圆柱几何校验。\\
\bottomrule
\end{tabularx}
\end{minipage}
\end{center}

\subsection{最优策略与时长增益}
最高精度复算得到
\begin{equation}
\theta=5.134472^\circ,\quad
v=139.988137\ \mathrm{m/s},\quad
t_0=0.814372\ \mathrm s,\quad
\delta=0.114129\ \mathrm s,
\label{eq:q2-param}
\end{equation}
故起爆时刻为 $\tau=0.928502\,\mathrm s$。投放点和起爆点见表\ref{tab:q2-strategy}。

\begin{table}[tbp]
\centering
\caption{问题二单机单弹最终策略}
\label{tab:q2-strategy}
\small
\renewcommand{\arraystretch}{1.20}
\setlength{\tabcolsep}{4.0pt}
\begin{tabularx}{\textwidth}{>{\centering\arraybackslash}p{0.10\textwidth} >{\centering\arraybackslash}p{0.12\textwidth} >{\centering\arraybackslash}p{0.12\textwidth} >{\centering\arraybackslash}p{0.12\textwidth} >{\centering\arraybackslash}X >{\centering\arraybackslash}X}
\toprule
航向/$^\circ$ & 速度/(m/s) & 投放/s & 延迟/s & 投放点 $(x,y,z)$/m & 起爆点 $(x,y,z)$/m \\
\midrule
5.134472 & 139.988137 & 0.814372 & 0.114129 &
\makecell{$(17913.545,$\\$10.202,1800.000)$} &
\makecell{$(17929.458,$\\$11.632,1799.936)$} \\
\bottomrule
\end{tabularx}
\end{table}

对应有效遮蔽区间为
\[
E_1=[0.928734,\ 5.516793]\ \mathrm s,
\]
故本问优化策略下
\begin{equation}
\boxed{H_1=4.588059082\ \mathrm s}.
\label{eq:q2-result}
\end{equation}
在相同完整遮蔽判据下，问题一固定策略的时长为 $1.391656494\,\mathrm s$，本问增加 $3.196402588\,\mathrm s$。图\ref{fig:q2-result}显示，起爆时刻与有效区间起点仅相差约 $2.32\times10^{-4}\,\mathrm s$，说明优化将云团几乎直接送入早期有利视线区域，而不是像固定策略那样等待相对几何关系逐渐改善。

\begin{figure}[tbp]
  \centering
  \includegraphics[width=0.92\textwidth]{F04_Q2最优区间与策略增益.pdf}
  \caption{问题二最优投放--起爆时序与相对固定策略的时长增益}
  \label{fig:q2-result}
\end{figure}

\subsection{搜索稳定性与数值验证}
图\ref{fig:q2-verify}将参数盆地、分辨率收敛和独立几何核复核放在同一图中。多次 DE 的高质量候选集中在航向约 $5^\circ$、投放时刻约 $0.7\,\mathrm s$ 的窄区域，双重退火给出的独立可行解约为 $4.568\,\mathrm s$，虽然略低于最终值，但仍进入同一高质量量级；这说明结论并非由单一初始化偶然产生。

\begin{figure}[tbp]
  \centering
  \includegraphics[width=0.92\textwidth]{F03_Q2候选盆地与稳定性.pdf}
  \caption{问题二候选盆地、数值分辨率与独立几何核复核}
  \label{fig:q2-verify}
\end{figure}

固定该策略后，中、高两级离散分别得到 $4.588226318\,\mathrm s$ 和 $4.588059082\,\mathrm s$，差值为 $1.672\times10^{-4}\,\mathrm s$；独立有限圆柱几何核的最大绝对距离差为 $2.32\times10^{-5}\,\mathrm m$。因此该策略在搜索机制、离散精度和几何实现三个层面均表现稳定。由于主搜索仍属于随机启发式，本文将其表述为\textbf{稳定高质量解}，不作理论全局最优声明。

\FloatBarrier

\section{问题三：单机三弹的联合时序优化}

\subsection{从单烟幕判据到三烟幕联合判据}
本问仍由 FY1 干扰 M1，因此问题一、二中的导弹运动、无人机运动、有限圆柱可见集和点到有限视线段距离均保持不变，只把单枚烟幕扩展为共享同一航迹的三枚烟幕。继承关系见表\ref{tab:q3-inherit}。

\begin{table}[tbp]
\centering
\caption{问题三相对问题二的继承--调整--输出关系}
\label{tab:q3-inherit}
\small
\begin{tabularx}{\textwidth}{p{0.25\textwidth} X X}
\toprule
\textbf{继承} & \textbf{调整} & \textbf{输出} \\
\midrule
$\vec m_1(t)$、$\vec f_1(t)$、$\mathcal T$、$\mathcal V_1(t)$、$d_{\rm seg}$ & 一条共享航迹上增加3组投放/起爆时序；同机相邻投放间隔不小于 $1\,\mathrm s$；每条视线可由不同烟幕承担遮蔽 & 三弹联合有效集合 $E_1$、时长 $H_1$ 与各弹独立作用区间 \\
\bottomrule
\end{tabularx}
\end{table}

按投放先后编号，定义八维决策向量
\begin{equation}
\vec x=(\theta,v,t_1,t_2,t_3,\delta_1,\delta_2,\delta_3)\in\mathbb R^8.
\label{eq:q3-x}
\end{equation}
其中三枚烟幕共享 $\theta,v$。给定 $\vec x$ 后，令 $\mathcal C(t;\vec x)$ 为时刻 $t$ 的有效烟幕中心集合，则问题一的单烟幕距离指标自然推广为
\begin{equation}
D_1(t;\vec x)=
\max_{\vec q\in\mathcal V_1(t)}
\min_{\vec c\in\mathcal C(t;\vec x)}
 d_{\rm seg}\!\left(\vec c,[\vec m_1(t),\vec q]\right).
\label{eq:q3-D}
\end{equation}
若 $\mathcal C(t;\vec x)=\varnothing$，约定 $D_1(t;\vec x)=+\infty$。式\eqref{eq:q3-D}的量词次序具有明确含义：\textbf{先对烟幕取最小值，允许不同烟幕分别遮挡不同视线；再对可见目标点取最大值，保证整个有限目标的每一条可见视线都被覆盖。}

由此定义
\begin{equation}
E_1(\vec x)=\{t:D_1(t;\vec x)\le10\},
\qquad
H_1(\vec x)=\mu(E_1(\vec x)).
\label{eq:q3-obj}
\end{equation}
完整八维模型写为
\begin{equation}
\left\{
\begin{aligned}
\max_{\vec x}\quad & H_1(\vec x),\\
\text{s.t.}\quad
&0\le\theta<2\pi,\qquad 70\le v\le140,\\
&0\le t_1,\\
&t_2-t_1\ge1,\qquad t_3-t_2\ge1,\\
&0\le\delta_k\le\sqrt{2h_1/g},\quad k=1,2,3,\\
&t_k+\delta_k\le T_1,\quad k=1,2,3.
\end{aligned}
\right.
\label{eq:q3-model}
\end{equation}

\subsection{多支路搜索与活跃边界复核}
问题二的优质单弹轨迹为三弹问题提供了有物理意义的初值，但八维联合模型可能存在完全不同的时序结构，因此不能只围绕问题二最优点做局部扩展。求解阶段设置三条完整八维搜索支路：一条以优质轨迹为暖启动，一条采用独立分层初值，一条从全域随机初值增强探索；三条支路始终使用同一 $H_1(\vec x)$ 评价函数。

最优候选还可能落在 $t_1=0$、$t_2-t_1=1$ 或 $\delta_k=0$ 等低维边界，而随机种群对这些边界的命中率较低。因此全局搜索后再主动固定疑似活跃约束，构造4--5维边界子问题进行独立复核。该步骤的目的不是重新定义模型，而是检验“边界最优”是否由模型结构真实支撑。

\steptext{1}{多源初始化}{同时构造问题二暖启动、独立分层样本和全域随机样本。}
\steptext{2}{三支路八维搜索}{在同一可行域内执行三条完整八维全局支路，统一评价 $H_1(\vec x)$。}
\steptext{3}{候选重排与盆地精修}{用同一中高精度评价器重排前列解，保留互异盆地并逐轮缩域。}
\steptext{4}{主动边界复核}{固定候选中疑似活跃的投放、间隔与起爆延迟约束，建立低维子问题独立搜索。}
\steptext{5}{严格认证}{执行投放间隔回代、三级数值复算、地面接触诊断和独立几何核校验。}

\subsection{最终策略与联合遮蔽结构}
最终共享航向和速度分别为
\[
\theta=5.134472^\circ,\qquad v=139.988137\ \mathrm{m/s},
\]
三枚烟幕的投放、起爆时序见表\ref{tab:q3}。

\begin{table}[tbp]
\centering
\caption{问题三 FY1 三枚烟幕的最终时序}
\label{tab:q3}
\small
\begin{tabular}{cccccc}
\toprule
弹号 & 投放时刻/s & 起爆延迟/s & 起爆时刻/s & 投放（起爆）点 $(x,y,z)$/m & 独立完整遮蔽/s \\
\midrule
1 & 0.000000 & 0.000000 & 0.000000 & $(17800.000,0.000,1800.000)$ & 2.556006 \\
2 & 1.000000 & 0.000000 & 1.000000 & $(17939.426,12.528,1800.000)$ & 4.225141 \\
3 & 2.000000 & 0.000000 & 2.000000 & $(18078.853,25.056,1800.000)$ & 0.000000 \\
\bottomrule
\end{tabular}
\end{table}

联合有效集合为
\[
E_1=[0.999999,\ 7.403160]\ \mathrm s,
\]
因此
\begin{equation}
\boxed{H_1=6.403160400\ \mathrm s}.
\label{eq:q3-result}
\end{equation}
三个投放时刻恰好为 $0,1,2\,\mathrm s$，首投非负约束和两条最小投放间隔均处于活跃边界；前两弹均立即起爆，说明该方案倾向于尽早建立前后衔接的两段完整遮蔽。

\subsection{第三弹零边际现象的消融解释}
第三弹独立完整遮蔽时长为0，本身并不能说明它对联合目标没有作用。为避免把“独立时长为0”误解成“联合贡献必为0”，进一步直接比较该策略下的时间集合：第二弹单独形成约 $[1.000000,5.225141]$，第一弹单独形成约 $[4.847154,7.403160]$，两者的并集已经等于最终联合集合 $[0.999999,7.403160]$（端点差来自数值边界定位精度）。因此，对\textbf{这一三弹方案}而言，移除第三弹不会缩短已形成的完整遮蔽区间，第三槽位的边际贡献确为0。

这一结论只说明当前方案存在三弹参数冗余，\textbf{不等价于证明“全局最优的两弹问题与三弹问题最优值必然相同”}。因此，本文只作局部饱和解释，不将其外推为两弹问题的全局最优结论。图\ref{fig:q3-saturation}同时给出独立/联合区间及代表性时序扰动：第三弹投放后移 $0.10\,\mathrm s$ 或延迟增加 $0.30\,\mathrm s$ 时，联合时长不变；而第二弹延迟增加 $0.30\,\mathrm s$ 会使目标减少约 $1.279\,\mathrm s$，表明前两弹是决定联合区间的关键时序变量。

\begin{figure}[tbp]
  \centering
  \includegraphics[width=0.94\textwidth]{F05_Q3联合区间与边际饱和.pdf}
  \caption{问题三独立/联合遮蔽区间与第三弹零边际平坦方向}
  \label{fig:q3-saturation}
\end{figure}

\subsection{边界与数值验证}
主动边界搜索中，固定“首投为0、首间隔为1、前两弹立即起爆”的低维子问题能够独立复现 $6.403160400\,\mathrm s$；而破坏关键间隔结构后，现有边界搜索记录中的最好值降至约 $5.920017\,\mathrm s$。这说明 $t_1=0$、$t_2-t_1=1$ 与前两弹立即起爆并非单次随机搜索造成的偶然贴边，而是当前策略的重要活跃结构。

三级严格复算得到约 $6.403162$、$6.403163$、$6.403160\,\mathrm s$，中、高两级差约 $3.05\times10^{-6}\,\mathrm s$；独立几何核最大绝对差约 $5.92\times10^{-6}\,\mathrm m$。因此，本问的主要结论可以概括为：\textbf{前两枚烟幕在活跃边界上形成连续时域衔接，第三槽位在当前方案附近呈现零边际冗余方向。}

\FloatBarrier

\section{问题四：三机单弹的时域接力协同优化}

\subsection{三平台十二维联合模型}
问题四仍只干扰 M1，但由 FY1--FY3 各投放一枚烟幕。相较问题三，有限圆柱几何核和多烟幕联合判据完全继承，变化在于三枚烟幕分别来自三条独立航迹，不再共享航向和速度。其继承--调整关系见表\ref{tab:q4-inherit}。

\begin{table}[tbp]
\centering
\caption{问题四相对问题三的继承--调整--输出关系}
\label{tab:q4-inherit}
\small
\begin{tabularx}{\textwidth}{p{0.25\textwidth} X X}
\toprule
\textbf{继承} & \textbf{调整} & \textbf{输出} \\
\midrule
$\mathcal T$、$\mathcal V_1(t)$、$d_{\rm seg}$、多烟幕 $\max$--$\min$ 完整遮蔽判据 & 三枚烟幕分别来自 FY1--FY3；三组航向、速度和投放/起爆时序相互独立；无同机投弹间隔耦合 & 三机联合有效集合 $E_1$、时长 $H_1$，以及各机独立区间与联合区间的集合关系 \\
\bottomrule
\end{tabularx}
\end{table}

对 FY$i$ 定义四维变量块
\begin{equation}
\vec x_i=(\theta_i,v_i,t_i,\delta_i),\qquad i=1,2,3,
\end{equation}
总决策向量为
\begin{equation}
\vec x=(\vec x_1,\vec x_2,\vec x_3)\in\mathbb R^{12}.
\label{eq:q4-x}
\end{equation}
时刻 $t$ 的有效烟幕中心集合仍记为 $\mathcal C(t;\vec x)$，故 $D_1(t;\vec x)$ 继续按式\eqref{eq:q3-D}计算，目标为
\[
H_1(\vec x)=\mu\{t:D_1(t;\vec x)\le10\}.
\]
完整模型为
\begin{equation}
\left\{
\begin{aligned}
\max_{\vec x}\quad &H_1(\vec x),\\
\text{s.t.}\quad
&0\le\theta_i<2\pi,\quad i=1,2,3,\\
&70\le v_i\le140,\quad i=1,2,3,\\
&0\le t_i\le T_1,\quad i=1,2,3,\\
&0\le\delta_i\le\sqrt{2h_i/g},\quad i=1,2,3,\\
&t_i+\delta_i\le T_1,\quad i=1,2,3.
\end{aligned}
\right.
\label{eq:q4-model}
\end{equation}

记三枚烟幕单独使用时的有效集合为 $E_{\rm FY1},E_{\rm FY2},E_{\rm FY3}$。由于增加烟幕只会使式\eqref{eq:q3-D}中的内层最小值不增，因此恒有
\begin{equation}
E_{\rm FY1}\cup E_{\rm FY2}\cup E_{\rm FY3}\subseteq E_1.
\label{eq:q4-union-relation}
\end{equation}
若最终为严格包含，则存在“单枚均不能完整遮蔽、但多枚空间互补后可完整遮蔽”的额外时间；若最终取等号，则协同主要表现为不同无人机在时间轴上的接力。该集合关系为本问结果机制的直接判据。

\subsection{结构化候选与变量块精修}
十二维完全随机初始化会产生大量无法形成有效遮蔽的个体。为提高搜索有效率，先分别构造 FY1--FY3 的单机候选库，再将不同无人机的优质候选交叉组合并混入随机样本，形成完整十二维种群。全局阶段采用多条差分进化支路以保留探索多样性；候选统一严格复算后保留多个互异盆地。

局部阶段利用 $\vec x=(\vec x_1,\vec x_2,\vec x_3)$ 的自然块结构：一次只开放一架无人机的四个变量，其余两架暂时固定，依次轮换 FY1--FY3；每次候选比较仍重新计算完整三机联合目标。空间结构基本稳定后，再单独细化投放和起爆时序。分块只改变局部搜索组织方式，不改变十二维模型与联合评价函数。

\steptext{1}{单机候选预筛}{分别获得 FY1--FY3 的单机可行优质参数，建立结构化联合初值。}
\steptext{2}{十二维联合搜索}{多支路搜索完整模型\eqref{eq:q4-model}，统一按三烟幕联合 $H_1$ 评价。}
\steptext{3}{多盆地保留}{严格重排候选，根据目标值和参数距离保留多个互异盆地。}
\steptext{4}{三机变量块轮换}{依次开放 $\vec x_1,\vec x_2,\vec x_3$ 进行局部精修，并穿插小尺度跨机扰动。}
\steptext{5}{时序抛光与认证}{细化投放/起爆时刻，最后执行三级复算、约束回代和独立几何核校验。}

\subsection{最终策略与时域接力机制}
最终三架无人机的策略见表\ref{tab:q4}。

\begin{table}[tbp]
\centering
\caption{问题四三架无人机最终策略}
\label{tab:q4}
\small
\begin{tabular}{lccccc}
\toprule
无人机 & 航向/$^\circ$ & 速度/(m/s) & 投放/s & 起爆延迟/s & 起爆/s \\
\midrule
FY1 & 5.062709 & 130.262587 & 0.872988 & 0.122159 & 0.995147 \\
FY2 & 280.289844 & 140.000000 & 3.968104 & 5.833163 & 9.801267 \\
FY3 & 75.593388 & 140.000000 & 21.268033 & 1.530087 & 22.798120 \\
\bottomrule
\end{tabular}
\end{table}

三架无人机单独形成的有效区间分别为
\begin{align*}
E_{\rm FY1}&=[1.000517,5.587754]\ \mathrm s,\\
E_{\rm FY2}&=[10.378120,14.318231]\ \mathrm s,\\
E_{\rm FY3}&=[23.530826,26.672740]\ \mathrm s.
\end{align*}
三段区间彼此不重叠，且最终联合结果恰好满足
\[
E_1=E_{\rm FY1}\cup E_{\rm FY2}\cup E_{\rm FY3}.
\]
因此式\eqref{eq:q4-union-relation}在最终方案处取等号，表明本问没有依赖“同一时刻多烟幕空间拼接”产生额外完整遮蔽，而是由三架无人机依靠不同初始位置形成前、中、后三段\textbf{时域接力}。联合时长为
\begin{equation}
\boxed{H_1=11.669261475\ \mathrm s}.
\label{eq:q4-result}
\end{equation}
图\ref{fig:q4-relay}以同一时间轴直接展示了三段独立区间和联合并集，避免仅从三张分离结果图推断协同机制。

\begin{figure}[tbp]
  \centering
  \includegraphics[width=0.94\textwidth]{F06_Q4时域接力与收敛.pdf}
  \caption{问题四三机时域接力结构与三级数值复算}
  \label{fig:q4-relay}
\end{figure}

\subsection{结果增益与数值验证}
现有求解记录显示，结构化联合候选已经达到约 $11.63\,\mathrm s$；完整十二维联合搜索进一步进入约 $11.677\,\mathrm s$ 的中精度候选层级，随后变量块精修和时序抛光只承担更细尺度的改善。最终严格 L1/L2/L3 复算得到 $11.675208$、$11.670319$ 和 $11.669261475\,\mathrm s$，中、高两级差约 $0.001057\,\mathrm s$；独立有限圆柱几何核的最大绝对距离差约 $7.27\times10^{-5}\,\mathrm m$。

FY2 与 FY3 的速度达到 $140\,\mathrm{m/s}$ 上界，而 FY1 约为 $130.26\,\mathrm{m/s}$，说明不同初始空间位置对应不同的最佳到达节奏。更重要的是，最终区间完全分离这一集合事实给出了比“参数达到上界”更直接的机制证据：\textbf{三机协同收益主要来自可达时间窗的纵向接续，而非同时多云团的横向空间互补。}

\FloatBarrier

\section{问题五：五机多弹对三导弹的四十维联合优化}

问题五同时继承问题三的“同机三弹共享航迹与投弹间隔”以及问题四的“多机独立航迹协同”，并把评价对象从单枚 M1 扩展为 M1--M3。为避免把后验覆盖关系误写成先验任务分配，正式模型始终让全部有效烟幕自由参与每枚导弹的完整遮蔽评价。继承--调整关系见表\ref{tab:q5-inherit}。

\begin{table}[tbp]
\centering
\caption{问题五相对问题三、问题四的继承--调整--输出关系}
\label{tab:q5-inherit}
\small
\begin{tabularx}{\textwidth}{p{0.27\textwidth} X X}
\toprule
\textbf{继承} & \textbf{调整} & \textbf{输出} \\
\midrule
有限圆柱可见集、多烟幕 $\max$--$\min$ 判据、同机投弹间隔、多机独立航迹 & 扩展为5架无人机、每机至多3枚烟幕和3枚导弹；目标由单个 $H_1$ 改为三导弹累计效能 $H_\Sigma$；不设置排他分配变量 & 40维联合策略、$E_1,E_2,E_3$、$H_1,H_2,H_3$、累计效能与墙钟重叠结构 \\
\bottomrule
\end{tabularx}
\end{table}

\subsection{三槽连续表示与原题资源约束的等价性}
对 FY$i$ 定义八维变量块
\begin{equation}
\vec x_i=(\theta_i,v_i,t_{i1},t_{i2},t_{i3},\delta_{i1},\delta_{i2},\delta_{i3})\in\mathbb R^8,
\qquad i=1,\ldots,5,
\label{eq:q5-block}
\end{equation}
总决策向量为
\begin{equation}
\vec x=(\vec x_1,\ldots,\vec x_5)\in\mathbb R^{40}.
\label{eq:q5-40d}
\end{equation}
其中 $t_{ik}$ 按投放先后排序，并满足
\begin{equation}
t_{i,k+1}-t_{ik}\ge1,
\qquad k=1,2.
\label{eq:q5-gap}
\end{equation}

题面规定“每架无人机至多投放3枚”。为了避免再引入5组离散启用变量，本文以每机3个连续槽位直接表示。该处理需要证明不会改变原题最优目标值。

\begin{propositionbox}
\textbf{命题2（三槽表示的目标等价性）}\quad
在假设3“不同烟幕之间不存在负向相互作用”成立、且本题任务时域允许补充满足 $1\,\mathrm s$ 间隔的投放时刻时，“每机至多3枚”的原问题与“每机固定3个连续槽位”的模型具有相同最优 $H_\Sigma$。

\textbf{证明：}
记原问题可行域为 $\Omega_{\le3}$，固定三槽可行域为 $\Omega_3$。显然 $\Omega_3\subseteq\Omega_{\le3}$，故三槽模型的最优值不超过原问题最优值。反之，任取 $\Omega_{\le3}$ 中的方案；若某架无人机实际少于3枚，则可在 $[0,T_{\max}]$ 内补充满足间隔约束的投放时刻并取 $\delta=0$。本题 $T_{\max}>60\,\mathrm s$，而每机只需布置至多3个相隔 $1\,\mathrm s$ 的时刻，因此补充时刻总能构造。加入额外烟幕后，有效烟幕集合只会扩大，式\eqref{eq:q5-D}中的内层最小距离不会增大，故各 $E_j$ 与 $H_j$ 均不会减小。于是任一“至多3枚”方案都可扩展为目标值不低于原方案的三槽方案，得到反向不等式。两者最优值相等。证毕。\proofend
\end{propositionbox}

因此，最终方案中出现边际贡献为0的槽位并非资源约束处理错误，而是连续等价表示允许存在的冗余槽位；其边际状态将在最终方案后单独核验。

\subsection{三导弹联合效能与完整模型}
给定40维策略 $\vec x$，令 $\mathcal C(t;\vec x)$ 为时刻 $t$ 仍处于有效期的全部烟幕中心集合。对导弹 $M_j$ 定义
\begin{equation}
D_j(t;\vec x)=
\max_{\vec q\in\mathcal V_j(t)}
\min_{\vec c\in\mathcal C(t;\vec x)}
 d_{\rm seg}\!\left(\vec c,[\vec m_j(t),\vec q]\right),
\qquad j=1,2,3,
\label{eq:q5-D}
\end{equation}
若 $\mathcal C(t;\vec x)=\varnothing$，约定 $D_j=+\infty$。于是
\begin{equation}
E_j(\vec x)=\{t:D_j(t;\vec x)\le10\},
\qquad
H_j(\vec x)=\mu(E_j(\vec x)).
\label{eq:q5-Hj}
\end{equation}
三枚导弹的主效能定义为
\begin{equation}
H_\Sigma(\vec x)=\sum_{j=1}^{3}H_j(\vec x),
\label{eq:q5-sum}
\end{equation}
单位记作“导弹$\cdot$s”，其中“导弹”为无量纲计数因子。完整模型写为
\begin{equation}
\left\{
\begin{aligned}
\max_{\vec x\in\mathbb R^{40}}\quad &H_\Sigma(\vec x),\\
\text{s.t.}\quad
&0\le\theta_i<2\pi,\qquad 70\le v_i\le140,\\
&0\le t_{i1},\qquad t_{i,k+1}-t_{ik}\ge1,\\
&0\le\delta_{ik}\le\sqrt{2h_i/g},\\
&t_{ik}+\delta_{ik}\le T_{\max},
\qquad i=1,\ldots,5,\ k=1,2,3.
\end{aligned}
\right.
\label{eq:q5-model}
\end{equation}
其中投放间隔式仅对 $k=1,2$ 使用。模型中不存在 UAV--导弹排他分配变量；同一烟幕可同时帮助多枚导弹，不同烟幕也可分别遮挡同一导弹视野中的不同目标射线。

为解释真实墙钟时间，另定义
\begin{equation}
H_\cup=\mu(E_1\cup E_2\cup E_3),
\label{eq:q5-union}
\end{equation}
但它只用于解释多导弹时间重叠，不替代主目标 $H_\Sigma$。

\subsection{结构化初始化、全局搜索与局部精修}
40维完全随机初始化在既有全局搜索中可用候选率较低。为此先求解 $5\times3=15$ 个“单无人机--单导弹”的三弹局部子问题，构造可达能力矩阵
\begin{equation}
A=(a_{ij})=
\begin{pmatrix}
4.9449&0&0\\
3.8473&3.7904&3.3741\\
3.0404&2.9475&3.0693\\
3.7890&3.7243&2.6882\\
3.8485&3.5867&3.4790
\end{pmatrix}\ \mathrm s.
\label{eq:q5-A}
\end{equation}
$a_{ij}$ 只描述 FY$i$ 对 $M_j$ 的局部可达性，用于生成互补的40维初值；进入正式搜索后不再出现在目标函数与约束中。

全局阶段设置4条完整40维 DE 支路：前三条的初始种群含结构化或局部能力块，第四条为纯随机冷启动。各支路在相同120代预算下产生候选后，统一用严格评价器复算，并按目标值与参数距离筛选互异盆地。局部阶段利用 $\vec x=(\vec x_1,\ldots,\vec x_5)$ 的自然块结构，依次开放一架无人机的8个变量进行块精修；空间结构稳定后，只对30个投放/起爆时序变量执行 $0.020$、$0.005$、$0.00125$ 和 $0.0003125\,\mathrm s$ 四级后抛光。整个过程始终评价同一个 $H_\Sigma$，分块仅改变搜索组织方式。

\steptext{1}{结构化种子构造}{由15个局部子问题与随机样本形成多样化40维初始种群。}
\steptext{2}{四支路完整搜索}{在同一40维模型上运行结构化、混合和纯随机冷启动支路。}
\steptext{3}{统一严格重排}{对前列候选重新计算 $(H_1,H_2,H_3,H_\Sigma)$ 并筛选互异盆地。}
\steptext{4}{五机变量块轮换}{依次开放五个8维变量块，所有接受判定仍基于完整40维联合目标。}
\steptext{5}{时序后抛光}{固定空间结构，以四级步长细化投放/起爆时序。}
\steptext{6}{严格认证}{执行 L1/L2/L3 分层复算、约束回代、地面诊断与独立有限圆柱几何复核。}

图\ref{fig:q5flow}只展示求解控制逻辑，与上述 STEP 分工：STEP 给出顺序与职责，流程图突出“盆地精修--收敛判断--时序抛光--认证”的循环关系。
\begin{figure}[H]
  \centering
  \includegraphics[width=0.98\textwidth]{F07_Q5联合优化流程图.pdf}
  \caption{问题五40维联合模型的求解与严格认证流程}
  \label{fig:q5flow}
\end{figure}

\subsection{最终方案与烟幕槽位作用}
五架无人机最终航向与速度见表\ref{tab:q5-uav}。
\begin{table}[H]
\centering
\caption{问题五五架无人机最终飞行策略}
\label{tab:q5-uav}
\begin{tabular}{lcc}
\toprule
无人机 & 航向/$^\circ$ & 速度/(m/s) \\
\midrule
FY1 & 179.465931 & 121.850174 \\
FY2 & 262.125542 & 117.012280 \\
FY3 & 85.612140  & 139.996881 \\
FY4 & 264.621692 & 139.527140 \\
FY5 & 115.183105 & 139.917936 \\
\bottomrule
\end{tabular}
\end{table}

十五个连续烟幕槽的投放--起爆时序见表\ref{tab:q5-time}。
\begin{table}[tbp]
\centering
\caption{问题五十五个烟幕槽的投放与起爆时序}
\label{tab:q5-time}
\small
\begin{tabular}{lccccc}
\toprule
无人机 & 槽号 & 投放时刻/s & 起爆延迟/s & 起爆时刻/s & 事后主要贡献标签 \\
\midrule
FY1&1&0.163282&3.412365&3.575648&M1\\
FY1&2&3.816973&4.948943&8.765917&M1\\
FY1&3&10.473130&14.342959&24.816089&M1\\
FY2&1&5.109965&6.538273&11.648238&M1\\
FY2&2&6.360779&12.025963&18.386742&M3\\
FY2&3&38.266413&13.190887&51.457300&M1\\
FY3&1&0.239507&6.550830&6.790337&M3\\
FY3&2&23.417780&0.130784&23.548564&M2\\
FY3&3&49.189090&2.993653&52.182743&M1\\
FY4&1&0.582916&10.977226&11.560142&M2\\
FY4&2&5.236528&11.567951&16.804479&M3\\
FY4&3&56.125653&3.543556&59.669209&M1\\
FY5&1&12.457206&0.262056&12.719262&M3\\
FY5&2&30.619797&5.880122&36.499920&M1\\
FY5&3&45.072508&9.069050&54.141558&M1\\
\bottomrule
\end{tabular}
\end{table}

图\ref{fig:q5deploy}把15个槽放在同一时间轴上，并用第二编码区分正边际与零边际槽位。
\begin{figure}[tbp]
  \centering
  \includegraphics[width=0.90\textwidth]{F08_Q5部署时序.pdf}
  \caption{问题五五机十五弹的投放--起爆时序与边际状态}
  \label{fig:q5deploy}
\end{figure}

逐弹删除审计显示，15个槽中7个对最终 $H_\Sigma$ 有正边际贡献，8个槽的边际贡献为0。图\ref{fig:q5marginal}进一步把每个槽对 M1--M3 的边际贡献展开：FY1-1、FY1-2、FY2-1 主要支撑 M1，FY3-2 与 FY4-1 主要支撑 M2，FY4-2 与 FY5-1 主要支撑 M3。零边际槽与命题2相一致，它们用于保持连续三槽表示的统一性，而不应被解释为“必须消耗才能达到当前效能”的关键烟幕。

\begin{figure}[tbp]
  \centering
  \includegraphics[width=0.78\textwidth]{F10_Q5逐弹边际贡献矩阵.pdf}
  \caption{问题五十五个烟幕槽对三枚导弹的逐弹边际贡献}
  \label{fig:q5marginal}
\end{figure}

\subsection{三导弹累计效能与重叠结构}
最终有效区间见表\ref{tab:q5-intervals}。
\begin{table}[tbp]
\centering
\caption{问题五三枚导弹的联合有效遮蔽区间}
\label{tab:q5-intervals}
\begin{tabular}{lccc}
\toprule
导弹 & 区间 & 起点/s & 终点/s \\
\midrule
M1&1&4.738335&10.950872\\
M1&2&12.377999&16.248967\\
M2&1&12.829424&16.545054\\
M2&2&24.263066&27.265874\\
M3&1&13.218021&17.086223\\
M3&2&17.209413&20.703855\\
\bottomrule
\end{tabular}
\end{table}

因此
\begin{equation}
H_1=10.083504639\,\mathrm s,\qquad
H_2=6.718437500\,\mathrm s,\qquad
H_3=7.362644043\,\mathrm s,
\label{eq:q5-H-result}
\end{equation}
并得到
\begin{equation}
\boxed{H_\Sigma=24.164586182\ \text{导弹}\cdot\mathrm s},
\qquad H_\cup=17.418010254\,\mathrm s.
\label{eq:q5-final}
\end{equation}

按同一墙钟时刻受干扰导弹数量分解，记恰有1枚、2枚、3枚导弹同时受干扰的墙钟时长为 $W_1,W_2,W_3$，最终结果为
\[
W_1=13.702380\,\mathrm s,\qquad
W_2=0.684684\,\mathrm s,\qquad
W_3=3.030946\,\mathrm s,
\]
故
\begin{equation}
H_\cup=W_1+W_2+W_3,
\qquad
H_\Sigma=W_1+2W_2+3W_3.
\label{eq:overlap-id}
\end{equation}
图\ref{fig:q5joint}显示，约 $13.218\sim16.249\,\mathrm s$ 内三枚导弹同时处于有效遮蔽状态，形成 $3.031\,\mathrm s$ 的三重重叠；另有 $0.685\,\mathrm s$ 的双重重叠。因而 $H_\Sigma>H_\cup$ 并非重复计数错误，而是主效能对同时干扰多枚导弹的真实累计。

\begin{figure}[tbp]
  \centering
  \includegraphics[width=0.88\textwidth]{F09_Q5联合区间与重叠.pdf}
  \caption{问题五三枚导弹的联合有效区间与同时受干扰数量}
  \label{fig:q5joint}
\end{figure}

\subsection{初始化消融、块精修与严格复算}
全局搜索记录可直接形成初始化方式的对照。前三条含结构化/混合初值的40维支路，在相同120代预算后经严格复算得到的支路最优候选分别为 $19.1803$、$18.9478$ 和 $18.9280\,\text{导弹}\cdot\mathrm s$；纯随机冷启动 DE4 的支路最优候选仅为 $1.4395\,\text{导弹}\cdot\mathrm s$。这说明在当前昂贵、非光滑的40维评价下，结构化初值明显提高了有限预算内进入高质量可行盆地的能力。该比较来自现有单次分支记录，因而只作为初始化机制的消融证据，不解释为跨算法统计显著性检验。

进入最佳盆地后，五机变量块轮换把严格候选从约 $18.93$ 继续提升至 $24.26\,\text{导弹}\cdot\mathrm s$ 量级，随后时序抛光只产生 $10^{-2}\sim10^{-4}\,\text{导弹}\cdot\mathrm s$ 的细尺度改进。图\ref{fig:q5verify}将初始化消融、最佳盆地逐轮抬升和最终 L1--L3 复算放在同一证据链中。

\begin{figure}[tbp]
  \centering
  \includegraphics[width=0.92\textwidth]{F11_Q5搜索改进与数值收敛.pdf}
  \caption{问题五结构化初始化、块精修与分层严格复算}
  \label{fig:q5verify}
\end{figure}

最终 L1/L2/L3 复算见表\ref{tab:q5-conv}。
\begin{table}[tbp]
\centering
\caption{问题五最终方案的分层数值复算}
\label{tab:q5-conv}
\small
\begin{tabular}{cccccc}
\toprule
层级 & 时间步长/s & 平面网格 & $H_1$/s & $H_2$/s & $H_3$/s \\
\midrule
L1 & 0.025 & 101 & 10.085388 & 6.723315 & 7.363660 \\
L2 & 0.015 & 151 & 10.083867 & 6.719531 & 7.362902 \\
L3 & 0.010 & 201 & 10.083504639 & 6.718437500 & 7.362644043 \\
\bottomrule
\end{tabular}
\end{table}
L2--L3 的最大分导弹时长差为 $0.001094\,\mathrm s$；独立有限圆柱几何核与高精度评价器之间的最大距离差为 $0.000367\,\mathrm m$。全部速度、投放间隔、起爆高度与任务时限约束均通过回代。因此，本问结果可表述为经“初始化消融--多盆地搜索--变量块精修--时序抛光--分层复算--独立几何核”支持的\textbf{稳定高质量方案}，但不作理论全局最优声明。

\FloatBarrier

\section{模型检验与评价}

\subsection{跨问题数值收敛与几何一致性}
五问均在最终策略上进行末两级分辨率比较，并用独立有限圆柱几何实现复核关键距离。图\ref{fig:validation}统一给出两类绝对误差，不再用人为门限比值替代原始数值。Q1--Q5 的末两级时长差分别约为 $6.10\times10^{-6}$、$1.67\times10^{-4}$、$3.05\times10^{-6}$、$1.06\times10^{-3}$ 和 $1.09\times10^{-3}\,\mathrm s$；独立几何核最大绝对差分别约为 $8.66\times10^{-5}$、$2.32\times10^{-5}$、$5.92\times10^{-6}$、$7.27\times10^{-5}$ 和 $3.67\times10^{-4}\,\mathrm m$。相对于 $10\,\mathrm m$ 的烟幕有效半径，最大的几何差约为 $3.7\times10^{-5}$ 的相对量级。

\begin{figure}[tbp]
  \centering
  \includegraphics[width=0.82\textwidth]{F12_全题统一验证.pdf}
  \caption{五个问题的末级数值收敛与独立有限圆柱几何复核}
  \label{fig:validation}
\end{figure}

这些结果说明：最终关键时长在继续加密时间/视线分辨率后只发生毫秒或更小变化，同时两套几何实现对关键距离的差异远小于烟幕半径。数值证据支持本文结果对离散尺度和几何实现的稳定性，但不等价于解析误差上界。

\subsection{模型优势及其证据}
\begin{enumerate}[label=\arabic*.]
  \item \textbf{几何对象与题意一致。} 问题一直接在有限闭圆柱的首次可见点集合上定义最坏视线距离，并用有限线段投影截断排除导弹后方/目标后方的伪相交；独立几何核误差最高仅 $3.67\times10^{-4}\,\mathrm m$，说明该几何口径在实现层面保持一致。
  \item \textbf{多烟幕联合判据保持正确量词顺序。} 式\eqref{eq:q3-D}与\eqref{eq:q5-D}先对有效烟幕取最小距离，再对目标可见点取最大值，使不同烟幕可以分别遮挡不同视线，同时仍要求整个有限目标被完整覆盖。
  \item \textbf{跨问扩展只增加必要自由度。} 从4维单弹、8维同机三弹、12维三机到40维五机多弹，公共运动与几何评价核不变；问题三、问题四分别通过“第三槽零边际”和“三段区间完全分离”给出了可解释的结构结论，而非只输出参数表。
  \item \textbf{搜索与认证职责分离。} 启发式算法只负责产生候选；论文结论统一来自严格复算、分辨率加密、约束回代和独立几何核。问题五又利用纯随机冷启动分支形成初始化消融，进一步说明结构化种子的实际作用。
  \item \textbf{问题五保留跨导弹真实协同。} 模型不设置排他任务分配，同一烟幕可以同时服务多枚导弹；式\eqref{eq:overlap-id}把双重、三重同时干扰与 $H_\Sigma-H_\cup$ 的差异直接对应起来。
\end{enumerate}

\subsection{局限、改进与推广边界}
本文按题面采用硬边界球形烟幕、恒定下沉速度和确定性运动学，未考虑风场、浓度衰减、云团形变、投放误差和导弹主动机动。若用于更真实任务，应把云团中心漂移、有效半径和起爆误差建模为随机或区间参数，并在完整遮蔽判据外层加入机会约束、场景鲁棒性或最坏情形评价；此时本文的有限视线段几何核仍可保留，但有效烟幕集合 $\mathcal C(t)$ 与阈值条件需要重构。

另一方面，问题五的目标评价开销高、目标函数非光滑且40维搜索仍依赖随机启发式。当前证据能够支持“稳定高质量方案”，不能提供理论全局最优保证。后续可利用遮蔽区间边界与局部可达能力构造代理模型或自适应采样\cite{ref8}，并在统一函数评价预算下进行更多独立随机种子和异质算法对照。模型可推广到“移动观察者--有限目标--时变遮挡域”的协同规划问题，但必须重新核验目标几何、遮挡域动力学和资源约束，不能直接照搬本题参数。

\AIUsageDeclaration{语言润色、数学表达检查、程序调试与排版辅助}

\begin{thebibliography}{99}
\bibitem{ref1} 全国大学生数学建模竞赛组织委员会. 2025高教社杯全国大学生数学建模竞赛A题：烟幕干扰弹的投放策略[Z]. 2025.
\bibitem{ref2} Storn R, Price K. Differential Evolution -- A Simple and Efficient Heuristic for Global Optimization over Continuous Spaces[J]. Journal of Global Optimization, 1997, 11(4): 341--359. DOI: \href{https://doi.org/10.1023/A:1008202821328}{10.1023/A:1008202821328}.
\bibitem{ref3} Price K V, Storn R M, Lampinen J A. Differential Evolution: A Practical Approach to Global Optimization[M]. Berlin: Springer, 2005. DOI: \href{https://doi.org/10.1007/3-540-31306-0}{10.1007/3-540-31306-0}.
\bibitem{ref4} Kirkpatrick S, Gelatt C D, Vecchi M P. Optimization by Simulated Annealing[J]. Science, 1983, 220(4598): 671--680. DOI: \href{https://doi.org/10.1126/science.220.4598.671}{10.1126/science.220.4598.671}.
\bibitem{ref5} Nocedal J, Wright S J. Numerical Optimization[M]. 2nd ed. New York: Springer, 2006. DOI: \href{https://doi.org/10.1007/978-0-387-40065-5}{10.1007/978-0-387-40065-5}.
\bibitem{ref6} Ericson C. Real-Time Collision Detection[M]. San Francisco: Morgan Kaufmann, 2004.
\bibitem{ref7} Virtanen P, Gommers R, Oliphant T E, et al. SciPy 1.0: Fundamental Algorithms for Scientific Computing in Python[J]. Nature Methods, 2020, 17: 261--272. DOI: \href{https://doi.org/10.1038/s41592-019-0686-2}{10.1038/s41592-019-0686-2}.
\bibitem{ref8} Jones D R, Schonlau M, Welch W J. Efficient Global Optimization of Expensive Black-Box Functions[J]. Journal of Global Optimization, 1998, 13: 455--492. DOI: \href{https://doi.org/10.1023/A:1008306431147}{10.1023/A:1008306431147}.
\bibitem{ref9} Das S, Suganthan P N. Differential Evolution: A Survey of the State-of-the-Art[J]. IEEE Transactions on Evolutionary Computation, 2011, 15(1): 4--31. DOI: \href{https://doi.org/10.1109/TEVC.2010.2059031}{10.1109/TEVC.2010.2059031}.
\end{thebibliography}

\appendix

\section{支撑材料与附件包说明}
随文交付的压缩包按照“论文--代码--结果--图片”四类组织，所有文件均与本文当前模型、数值结果和图表一一对应，仅保留论文复现与提交所需材料。各目录用途如下：论文目录保存无代码版、含代码版及可直接编译的 \LaTeX{} 源文件；代码目录保存问题一至问题五的完整 Python 程序；结果目录按五个问题分别保存正式 CSV/JSON 结果、数值收敛记录、几何复核及 result1--result3 填表数据；图片目录保存统一重构的13幅正式图（F00--F12），均同时提供 PDF 与 PNG；绘图代码目录保存统一重绘脚本，所有结果图只读取正式 CSV/JSON 或固定参数做确定性后处理。

\begin{longtable}{p{0.08\textwidth} p{0.32\textwidth} p{0.50\textwidth}}
\toprule
序号 & 目录或文件 & 内容及与正文的对应关系 \\
\midrule
\endfirsthead
\toprule
序号 & 目录或文件 & 内容及与正文的对应关系 \\
\midrule
\endhead
1 & \texttt{论文/论文.pdf} & 无完整代码附录的正式论文版本，包含正文、参考文献、支撑材料说明及关键结果表。 \\
2 & \texttt{论文/论文\_含代码.pdf} & 在同一论文正文后追加问题一至问题五完整 Python 源程序，便于完整查阅。 \\
3 & \texttt{论文/main.tex}、\texttt{main\_含代码.tex} & 两个 PDF 对应的 \LaTeX{} 源文件；\texttt{cumcmthesis.cls} 为排版类文件。 \\
4 & \texttt{代码/1.py--5.py} & 问题一至问题五完整求解程序，分别对应正文第6--10章。 \\
5 & \texttt{结果/问题1--问题5/} & 各问正式 CSV/JSON 结果、数值收敛、几何复核和 result1--result3 对应数据。 \\
6 & \texttt{图片/重构图/PDF/} & F00--F12 共13幅正式论文图的 PDF 矢量/高质量原图，包括技术路线、几何机理、结果与验证图。 \\
7 & \texttt{图片/重构图/PNG/} & 与 F00--F12 逐一对应的 PNG 预览图，便于审图与其他文档引用。 \\
8 & \texttt{绘图代码/论文图重绘.py} & 统一重绘脚本；只读取正式结果或固定参数做确定性后处理，不承担优化求解。 \\
\bottomrule
\end{longtable}

\section{结果表关键坐标}

\subsection{问题三 result1 对应数据}
\begin{table}[H]
\centering
\caption{问题三三枚烟幕弹投放与起爆坐标}
\small
\begin{tabular}{cccc}
\toprule
弹号 & 投放点 $(x,y,z)$/m & 起爆点 $(x,y,z)$/m & 联合方案说明 \\
\midrule
1 & $(17800.000,0.000,1800.000)$ & $(17800.000,0.000,1800.000)$ & 共享航迹 \\
2 & $(17939.426,12.528,1800.000)$ & $(17939.426,12.528,1800.000)$ & 共享航迹 \\
3 & $(18078.853,25.056,1800.000)$ & $(18078.853,25.056,1800.000)$ & 共享航迹 \\
\bottomrule
\end{tabular}
\end{table}

\subsection{问题四 result2 对应数据}
\begin{table}[H]
\centering
\caption{问题四三枚烟幕弹投放与起爆坐标}
\scriptsize
\begin{tabular}{lcc}
\toprule
无人机 & 投放点 $(x,y,z)$/m & 起爆点 $(x,y,z)$/m \\
\midrule
FY1 & $(17913.274048,10.035131,1800.000000)$ & $(17929.124678,11.439364,1799.926879)$ \\
FY2 & $(12099.233931,853.400236,1400.000000)$ & $(12245.109057,49.891715,1233.273626)$ \\
FY3 & $(6740.813095,-116.105256,700.000000)$ & $(6794.109442,91.370904,688.528286)$ \\
\bottomrule
\end{tabular}
\end{table}

\subsection{问题五 result3 对应数据}
\begin{table}[H]
\centering
\caption{问题五 result3 关键坐标与事后主要贡献标签}
\tiny
\setlength{\tabcolsep}{2.4pt}
\resizebox{0.99\textwidth}{!}{%
\begin{tabular}{llrrrrrrl}
\toprule
无人机 & 弹号 & 投放$x$/m & 投放$y$/m & 起爆$x$/m & 起爆$y$/m & 起爆$z$/m & 独立时长/s & 标签 \\
\midrule
FY1&1&17780.105&0.185&17364.326&4.061&1742.943&4.076458&M1\\
FY1&2&17334.921&4.335&16731.918&9.956&1679.989&2.184954&M1\\
FY1&3&16523.903&11.895&14776.287&28.186&791.970&0.000000&M1\\
FY2&1&11918.082&807.709&11813.267&49.865&1190.530&3.870968&M1\\
FY2&2&11898.030&662.729&11705.242&-731.188&691.343&0.000000&M3\\
FY2&3&11386.550&-3035.419&11175.086&-4564.361&547.402&0.000000&M1\\
FY3&1&6002.565&-2966.568&6072.730&-2052.160&489.724&0.000000&M3\\
FY3&2&6250.824&268.807&6252.225&287.063&699.916&3.002808&M2\\
FY3&3&6526.857&3866.135&6558.921&4284.009&656.086&0.000000&M1\\
FY4&1&10992.377&1919.025&10848.816&394.147&1209.553&3.715630&M2\\
FY4&2&10931.516&1272.579&10780.230&-334.358&1144.294&3.494441&M3\\
FY4&3&10265.985&-5796.576&10219.642&-6288.821&1738.472&0.000000&M1\\
FY5&1&12258.337&-422.680&12242.736&-389.499&1299.664&3.740276&M3\\
FY5&2&11176.994&1877.051&10826.911&2621.587&1130.578&0.000000&M1\\
FY5&3&10316.526&3707.040&9776.584&4855.355&896.986&0.000000&M1\\
\bottomrule
\end{tabular}%
}
\end{table}
\noindent\textit{注：}表中“标签”为最终联合解的事后主要贡献说明，不进入无人机--导弹联合优化的决策变量与约束。完整三维投放坐标、投放时刻、起爆延迟与起爆时刻见附件包 \texttt{结果/问题5/问题5\_result3填表数据.csv}。

\end{document}
